\DocumentMetadata{
  pdfversion=2.0,
  pdfstandard=ua-2,
  lang=en-US,
  tagging=on,
  tagging-setup={math/setup=mathml-SE,tabsorder=structure}
}
\documentclass[11pt,letterpaper]{article}
\usepackage[margin=0.68in,includefoot]{geometry}
\usepackage{newcomputermodern}
\usepackage{amsmath,amscd,mathtools}
\usepackage{tikz,adjustbox}
\providecommand{\square}{\mdlgwhtsquare}
\usepackage[hidelinks]{hyperref}
\hypersetup{
  pdftitle={Numerical Analysis First-Year Exam, August 2024},
  pdfauthor={University of Florida Department of Mathematics},
  pdfsubject={Graduate mathematics examination},
  pdfdisplaydoctitle=true,
  bookmarksopen=true
}
\setcounter{secnumdepth}{0}
\setlength{\parindent}{0pt}
\setlength{\parskip}{0.28em}
\setlength{\emergencystretch}{3em}
\setlength{\leftmargini}{2.15em}
\setlength{\leftmarginii}{2.65em}
\setlength{\leftmarginiii}{3em}
\pagestyle{empty}
\raggedbottom
\allowdisplaybreaks
\NewTaggingSocketPlug{block/list/label}{examproblem}{%
  \tagstructbegin{tag=Lbl}\tagstructend%
  \tagstructbegin{tag=\LBody}%
  \tagstructbegin{tag=\ProblemHeadingTag,title-o={\ProblemHeadingText}}%
  \tagmcbegin{tag=\ProblemHeadingTag,actualtext={\ProblemHeadingText}}%
  #2%
  \tagmcend%
  \tagstructend%
}
\newenvironment{examproblems}[2]{%
  \def\ProblemHeadingTag{#2}%
  \AssignTaggingSocketPlug{block/list/label}{examproblem}%
  \begin{enumerate}%
  \setcounter{enumi}{#1}%
  \renewcommand{\labelenumi}{\textbf{\theenumi.}}%
  \setlength{\topsep}{0.35em}%
  \setlength{\partopsep}{0pt}%
  \setlength{\itemsep}{0.48em}%
  \setlength{\parsep}{0.18em}%
}{\end{enumerate}}
\newenvironment{examsubparts}{%
  \AssignTaggingSocketPlug{block/list/label}{default}%
  \begin{enumerate}%
  \setlength{\topsep}{0.18em}%
  \setlength{\partopsep}{0pt}%
  \setlength{\itemsep}{0.16em}%
  \setlength{\parsep}{0.1em}%
}{\end{enumerate}}
\newenvironment{examinstructions}{%
  \par\begingroup\itshape
}{%
  \par\endgroup\vspace{0.18em}
}
\makeatletter
\renewcommand\subsection{\@startsection{subsection}{2}{0pt}%
  {-1.25ex plus -.25ex minus -.1ex}%
  {0.55ex plus .1ex}%
  {\normalfont\large\bfseries}}
\renewcommand{\@maketitle}{%
  \newpage\null\vskip -1em%
  {\footnotesize\bfseries
    \href{https://www.ufl.edu/}{University of Florida}, \href{https://math.ufl.edu/}{Department of Mathematics}\par}%
  \vskip -0.5em%
  \tagstructbegin{tag=H1,title={Numerical Analysis First-Year Exam, August 2024}}%
  \tagpdfparaOff%
  \tagmcbegin{tag=H1}%
    {\LARGE\bfseries\href{https://gma.math.ufl.edu/exams/numerical-analysis-first-year/}{Numerical Analysis First-Year Exam}, August 2024\par}%
  \tagmcend%
  \tagpdfparaOn%
  \tagstructend%
  \vskip 0.25em%
  \tagmcbegin{artifact}%
    \hrule height 1pt%
  \tagmcend%
  \vskip 1em%
}
\makeatother
\title{Numerical Analysis First-Year Exam, August 2024}
\author{}
\date{}
\begin{document}
\maketitle
\thispagestyle{empty}
\begin{examinstructions}
Do 4 (four) problems.
\end{examinstructions}
\begin{examproblems}{0}{H2}
\def\ProblemHeadingText{Problem 1.}
\item
Consider an approximate integration scheme of the following form, for \(0<\alpha<1\),
\[
\int_0^1 x^\alpha f(x)\,dx \approx A\int_0^1 f(x)\,dx+B\int_0^1 xf(x)\,dx.
\]
\begin{examsubparts}
\item[\textnormal{(a)}]
Determine the constants \(A,B\) so that the formula has the maximum degree of exactness.
\item[\textnormal{(b)}]
What is the degree of exactness of the formula in (a)?
\end{examsubparts}
\def\ProblemHeadingText{Problem 2.}
\item
Let~\(f\) be an arbitrary (continuous) function on \([0,1]\) satisfying \(f(x)+f(1-x)\equiv 1\) for \(0\leq x\leq 1\).
\begin{examsubparts}
\item[\textnormal{(a)}]
Show that \(\int_0^1 f(x)\,dx=\frac12\).
\item[\textnormal{(b)}]
Show that the composite trapezoidal rule for computing \(\int_0^1 f(x)\,dx\) is exact.
\end{examsubparts}
\def\ProblemHeadingText{Problem 3.}
\item
Discuss uniqueness and non-uniqueness of the least squares approximation to a function~\(f\) in the case of a discrete set \(T=\{t_1,t_2\}\) (i.e., \(N=2\)) where \(t_1\neq t_2\) and \(\Phi_n=P_{n-1}\) (polynomial of degree \(\leq n-1\)). In case of non-uniqueness, determine all solutions.
\def\ProblemHeadingText{Problem 4.}
\item
Derive the three-point formula for the second derivative
\[
f''(x_0)=\frac{1}{h^2}\bigl(f(x_0-h)-2f(x_0)+f(x_0+h)\bigr)-\frac{h^2}{12}f^{(4)}(\eta),
\]
 where~\(\eta\) is between \(x_0-h\) and \(x_0+h\).
\def\ProblemHeadingText{Problem 5.}
\item
Let \(x_0,x_1,\ldots,x_n\) be pairwise distinct points in \([a,b]\), \(-\infty<a<b<\infty\), and \(f\in C^1[a,b]\). Show that given any \(\epsilon>0\), there exists a polynomial~\(p\) such that \(\lVert f-p\rVert_\infty<\epsilon\) and at the same time \(p(x_i)=f(x_i)\), \(i=0,1,\ldots,n\). Here \(\lVert u\rVert_\infty=\max_{a\leq x\leq b}|u(x)|\).
\end{examproblems}
\end{document}
